\section{Project Understanding Model (PUM) Specification}

\subsection{Document Information}

\begin{table}[H]
\centering
\begin{tabular}{ll}
\toprule
Document & Project Understanding Model (PUM) Specification \\
Version & 1.0 \\
Status & Draft \\
Last Updated & July 2026 \\
Related Architecture & Engineering Shadow Reference Model v1 \\
\bottomrule
\end{tabular}
\end{table}

%%%%%%%%%%%%%%%%%%%%%%%%%%%%%%%%%%%%%%%%%%%%%%%%%%%%%%%%%%%%%%

\subsection{Purpose}

The Project Understanding Model (PUM) is the central knowledge component of the Engineering Shadow platform.

Its purpose is to continuously maintain an evolving understanding of an engineering project by integrating engineering evidence collected throughout the project lifecycle.

Rather than acting as a storage system for engineering artifacts, the PUM represents the Shadow's current understanding of the project, enabling context-aware reasoning and intelligent engineering assistance.

Within the Engineering Shadow Reference Model (ESRM), the PUM serves as the primary output of the Understanding Engine and the primary input to the Assistance Engine.

%%%%%%%%%%%%%%%%%%%%%%%%%%%%%%%%%%%%%%%%%%%%%%%%%%%%%%%%%%%%%%

\subsection{Motivation}

Modern engineering projects generate large amounts of information including source code, documentation, simulation results, design reviews, research papers, engineering discussions and tool outputs.

Although these artifacts contain valuable information, they rarely capture the reasoning, relationships and project context required to understand how a project evolved.

Consequently, engineers repeatedly reconstruct project context before making decisions, reviewing designs or continuing previous work.

The Project Understanding Model addresses this problem by continuously transforming engineering evidence into an evolving project understanding that can be reused throughout the project lifecycle.

%%%%%%%%%%%%%%%%%%%%%%%%%%%%%%%%%%%%%%%%%%%%%%%%%%%%%%%%%%%%%%

\subsection{Definition}

The Project Understanding Model (PUM) is defined as:

\begin{quote}
\textit{
A continuously evolving representation of both the structural knowledge and the current situational state of an engineering project, enabling Engineering Shadow to reason about what the project is, how it evolved and what should happen next.
}
\end{quote}

The PUM does not attempt to replicate the complete cognitive state of an engineer.

Instead, it maintains sufficient project understanding to provide meaningful, context-aware engineering assistance.

%%%%%%%%%%%%%%%%%%%%%%%%%%%%%%%%%%%%%%%%%%%%%%%%%%%%%%%%%%%%%%

\subsection{Design Philosophy}

The design of the PUM follows four fundamental principles.

\begin{enumerate}

\item \textbf{Understanding over Storage}

The objective of the PUM is to maintain engineering understanding rather than archive engineering information.

\item \textbf{Relationships over Isolated Facts}

Engineering understanding emerges from relationships between project concepts rather than independent pieces of information.

\item \textbf{Continuous Evolution}

Project understanding evolves continuously as new engineering evidence becomes available.

\item \textbf{Reasoning before Assistance}

Engineering assistance is generated from project understanding rather than directly from raw engineering information.

\end{enumerate}

%%%%%%%%%%%%%%%%%%%%%%%%%%%%%%%%%%%%%%%%%%%%%%%%%%%%%%%%%%%%%%

\subsection{Core Responsibilities}

The Project Understanding Model is responsible for:

\begin{itemize}

\item Maintaining an evolving understanding of the engineering project.

\item Preserving engineering rationale throughout the project lifecycle.

\item Capturing relationships between engineering concepts.

\item Maintaining awareness of the project's current situation.

\item Providing a consistent project context for intelligent engineering assistance.

\item Supporting traceability between engineering evidence and project understanding.

\end{itemize}

%%%%%%%%%%%%%%%%%%%%%%%%%%%%%%%%%%%%%%%%%%%%%%%%%%%%%%%%%%%%%%

\subsection{Understanding Dimensions}

The PUM maintains understanding across multiple complementary dimensions.

\begin{enumerate}

\item Project Context

Overall understanding of the project's objectives, scope and engineering intent.

\item Architectural Understanding

Understanding of the system architecture and relationships between engineering components.

\item Decision Understanding

Understanding of why engineering decisions were made and how they influenced the project.

\item Requirement Understanding

Understanding of the engineering requirements, constraints and design drivers.

\item Verification Understanding

Understanding of verification activities and confidence in engineering outcomes.

\item Progress Understanding

Understanding of the current state of the project, ongoing work and active engineering priorities.

\item Issue Understanding

Understanding of unresolved engineering problems, blockers and associated context.

\end{enumerate}

%%%%%%%%%%%%%%%%%%%%%%%%%%%%%%%%%%%%%%%%%%%%%%%%%%%%%%%%%%%%%%

\subsection{Expected Capabilities}

A mature Project Understanding Model should enable Engineering Shadow to answer questions such as:

\begin{itemize}

\item What is this project trying to achieve?

\item What is the current engineering objective?

\item Why was a particular design decision made?

\item Which requirements influenced this architecture?

\item How has the architecture evolved over time?

\item What engineering issues remain unresolved?

\item What evidence supports the current design?

\item What should be considered before making the next engineering decision?

\end{itemize}

These questions define the expected behavior of the PUM rather than its internal implementation.

%%%%%%%%%%%%%%%%%%%%%%%%%%%%%%%%%%%%%%%%%%%%%%%%%%%%%%%%%%%%%%

\subsection{Interaction within the ESRM}

Within the Engineering Shadow Reference Model:

\begin{itemize}

\item The Information Acquisition Layer provides raw engineering information.

\item The Evidence Construction Layer transforms observations into engineering evidence.

\item The Understanding Engine updates the Project Understanding Model using the available evidence.

\item The Assistance Engine generates engineering assistance using the current Project Understanding Model.

\end{itemize}

Consequently, the PUM acts as the central knowledge interface between engineering understanding and engineering assistance.

%%%%%%%%%%%%%%%%%%%%%%%%%%%%%%%%%%%%%%%%%%%%%%%%%%%%%%%%%%%%%%

\subsection{Architectural Boundaries}

The Project Understanding Model intentionally does not define:

\begin{itemize}

\item Internal data structures.

\item Storage technologies.

\item Database schemas.

\item Knowledge graph implementations.

\item AI models.

\item Algorithms for evidence processing.

\item User interfaces.

\end{itemize}

These concerns belong to subsequent architectural specifications.

%%%%%%%%%%%%%%%%%%%%%%%%%%%%%%%%%%%%%%%%%%%%%%%%%%%%%%%%%%%%%%

\subsection{Future Architectural Work}

This specification establishes the conceptual role of the PUM.

Future documents will describe:

\begin{itemize}

\item Engineering Ontology

\item Engineering Entities

\item Engineering Relationships

\item Decision Episode Integration

\item Project Understanding Evolution

\item Evidence Integration

\item Internal PUM Architecture

\item Reasoning Mechanisms

\end{itemize}

%%%%%%%%%%%%%%%%%%%%%%%%%%%%%%%%%%%%%%%%%%%%%%%%%%%%%%%%%%%%%%

\subsection{Summary}

The Project Understanding Model is the central knowledge asset of Engineering Shadow.

Rather than functioning as a repository of engineering information, it maintains a continuously evolving understanding of the engineering project.

By separating project understanding from engineering assistance, the PUM provides a stable architectural foundation upon which future reasoning, automation and intelligent engineering capabilities can be built.